% INF-221 Tarea 1 2026-2 | Tomás San Martín | Rol 202473565-9
El comportamiento práctico reproduce las cotas teóricas en el régimen
asintótico, pero la constante oculta, el tipo de entrada y las decisiones de
implementación (mediana de tres, fusión con min-heap, umbral de corte en
Strassen) determinan qué algoritmo conviene en cada rango de $n$. El análisis
experimental es, por tanto, el complemento necesario de la notación asintótica
para seleccionar algoritmos en la práctica, ya que también entran en juego la jerarquía de memoria y la eficiencia de caché 
, que no se reflejan en la complejidad teórica, y que afectan en lo que podriamos llegar a esperar.